% ============================================================
%  LAB / LEARNING JOURNAL
%  Build with latexmk (LaTeX Workshop's default recipe).
%  Save the file and a PDF should appear; or press Ctrl+Alt+B.
% ============================================================
\documentclass[11pt]{article}

% --- Page geometry: comfortable margins for reading/note-taking ---
\usepackage[a4paper, margin=1in]{geometry}   % use letterpaper if you're in the US

% --- Fonts / encoding (modern, sensible defaults) ---
\usepackage[T1]{fontenc}
\usepackage{lmodern}

% --- Math: amsmath is the workhorse; amssymb/amsthm add symbols & theorem tools ---
\usepackage{amsmath, amssymb, amsthm}

% --- Units: siunitx handles numbers + units correctly, e.g. \SI{9.81}{m/s^2} ---
\usepackage{siunitx}

% --- Figures (the actual \includegraphics example below is commented out) ---
\usepackage{graphicx}

% --- Code listings: no external tools needed (unlike minted, which wants shell-escape) ---
\usepackage{listings}
\usepackage{xcolor}
\lstset{
  basicstyle=\ttfamily\small,
  keywordstyle=\color{blue!70!black},
  commentstyle=\color{green!50!black},
  stringstyle=\color{purple!70!black},
  numbers=left, numberstyle=\tiny\color{gray},
  breaklines=true, frame=single, language=Python
}

% --- hyperref makes the table of contents and references clickable. Load it LAST. ---
\usepackage[colorlinks=true, linkcolor=blue!50!black, urlcolor=blue!50!black]{hyperref}

% ------------------------------------------------------------
%  Custom tools
% ------------------------------------------------------------

% An "entry" is just a section with a date built into the heading,
% so each one shows up in the table of contents automatically.
\newcommand{\entry}[2]{\section[#2]{#1\quad\textnormal{\textit{#2}}}}
%   usage: \entry{2026-06-25}{What I learned about limits}

% A box for results worth flagging. Definitions/claims you want to find later.
\theoremstyle{definition}
\newtheorem{keyresult}{Key result}

% ------------------------------------------------------------
\title{Learning Journal}
\author{Ian Silber}                      % put your name here if you like
\date{\today}                  % the cover date; entries carry their own dates

\begin{document}
\maketitle
\tableofcontents               % grows automatically as you add entries
\vspace{1em}\hrule\vspace{1em}

% ============================================================
%  ENTRIES  —  copy the block below to start a new one
% ============================================================

\entry{2026-06-25}{Catching up: What I have learned during my first month}

A short paragraph on what you worked on. Inline math goes between dollar
signs, like $f(x) = x^2$. Displayed equations get their own line:
\begin{equation}
  \lim_{h \to 0} \frac{f(x+h) - f(x)}{h} = f'(x).
  \label{eq:derivative}        % label lets you reference it later with \eqref
\end{equation}
You can point back to Equation~\eqref{eq:derivative} from anywhere.

A physical quantity with units, typeset properly:
the acceleration due to gravity is about \SI{9.81}{\meter\per\second\squared}.

Flag something you want to find again:
\begin{keyresult}
  The derivative is the limit of the difference quotient, when that limit
  exists. ``Differentiable'' means this limit exists at the point in question.
\end{keyresult}

A few things I'm still unsure about:
\begin{itemize}
  \item one open question
  \item another thing to revisit
\end{itemize}

A bit of code I was working through:
\begin{lstlisting}
import numpy as np

def difference_quotient(f, x, h=1e-6):
    return (f(x + h) - f(x)) / h
\end{lstlisting}

% To add a figure later, drop the image file next to this .tex and uncomment:
% \begin{figure}[htbp]
%   \centering
%   \includegraphics[width=0.6\linewidth]{my-plot.png}
%   \caption{What the plot shows.}
%   \label{fig:my-plot}
% \end{figure}

% ------------------------------------------------------------
\entry{2026-06-26}{Next entry}

Start writing here.

\end{document}